

%% 
%% This is a LaTeX document generated by automatic conversion of a
%% Mathematica notebook using Mathematica 4.0.
%% 
%% This document uses special macros that are defined in a LaTeX 
%% package file with a name of the form notebook*.sty.  Appropriate 
%% commands for loading the package have been included in this document.
%% The LaTeX package files can be found in a directory with path:
%% 
%% /usr/local/mathematica/SystemFiles/IncludeFiles/TeX/
%% 
%% The LaTeX package used in this document may require that your 
%% LaTeX implementation support fonts other than the standard Computer 
%% Modern fonts that are used by LaTeX.  See the Additional Information 
%% section under the TeXSave entry in the Mathematica Reference Guide 
%% for more information.
%% 
%% Implementation-specific instructions for installing TeX-related files
%% can be found at this URL.
%% 
%% http://www.wolfram.com/support/FrontEnds/Export/TeX
%% 
%% 

\documentclass{article}
\usepackage{notebook2e, latexsym}


\begin{document}

\Title{SAMPLING CONTINUOUS-TIME SIGNALS}
\Subtitle{LECTURE 23}
\Subsubtitle{ELE314: LINEAR SIGNALS AND SYSTEMS}
\Subsubtitle{Mariusz Jankowski \copyright{}

University of Southern Maine

mjkcc@usm.maine.edu}
\dispSFinmath{<<\Mvariable{SignalPlot`}}
\begin{CellGroup}
\Section*{1 Introduction}


In this document we will explore the operations of sampling and interpolation, two operation that bridge continuous and discrete time domains. 


\end{CellGroup}
\begin{CellGroup}
\Section*{2 Experiments in sampling continuous-time signals}


Sampling is an operation that leads to a representation of a continuous-time signal by evenly spaced samples. It may be surprising to some that such
an operation produces meaningful results, but this phenomenon is very common in our everyday life. In this section, sampling will be analyzed from
a purely time-domain perspective. Consider first a simple sinusoidal signal and a sampling interval T.

\dispSFinmath{f[\Mvariable{t\_}]:=\cos [2\multsp \pi \multsp t]}
\begin{CellGroup}
\dispSFinmath{\Mfunction{Plot}[f[t],\{t,0,5\},\Mvariable{PlotLabel}\rightarrow \Mvariable{Se\NTilde al\multsp de\multsp tiempo\multsp continuo}];}
\mathGraphic{lecture23.tex_gr1.eps}

\end{CellGroup}
\begin{CellGroup}
\dispSFinmath{\MathBegin{MathArray}{l}
T=0.1;  \\
\noalign{\vspace{0.9375ex}}  \\
\MathBegin{MathArray}{l}
\Muserfunction{LinePlot}\big[\Mfunction{Table}\big[\{k\multsp T,f[k\multsp T]\},\big\{k,0,\frac{5}{T}\big\}\big],  \\
\noalign{\vspace{1.07292ex}}
\hspace{1.em} \Mvariable{PlotLabel}\rightarrow \Mvariable{Se\NTilde al\multsp muestreada}\big];\\
\MathEnd{MathArray}
\MathEnd{MathArray}}
\mathGraphic{lecture23.tex_gr2.eps}

\end{CellGroup}
Qualitatively, the samples give a fair representation of a sinusoidal signal. Now, consider a longer sampling interval. 

\begin{CellGroup}
\dispSFinmath{\MathBegin{MathArray}{l}
T=0.35;  \\
\noalign{\vspace{0.9375ex}}  \\
\MathBegin{MathArray}{l}
\Muserfunction{LinePlot}\big[\Mfunction{Table}\big[\{k\multsp T,f[k\multsp T]\},\big\{k,0,\frac{5}{T}\big\}\big],  \\
\noalign{\vspace{1.07292ex}}
\hspace{1.em} \Mvariable{PlotLabel}\rightarrow \Mvariable{Se\NTilde al\multsp muestreada},\Mvariable{PlotRange}->\Mvariable{All}\big];\\
\MathEnd{MathArray}
\MathEnd{MathArray}}
\mathGraphic{lecture23.tex_gr3.eps}

\end{CellGroup}
The resemblence to \inlineTFinmath{f(t)} has dissapeared. However, such superficial analysis may lead to incorrect conclusions about sampling rates.
The sampling rate is dependent on the frequency content of a signal. The theoretical result  will be developed in the next section.



\begin{CellGroup}
\Subsubsection*{Problem 1}
Consider sampling a signal with linearly increasing frequency, called a chirp signal. A fixed sampling rate will eventually result in an undersampled
signal. In this problem Play two chirp signals, sampled at 8192 and 5000 samples/sec. Listen for any differences in the two signals.

\begin{CellGroup}
\dispSFinmath{\Mfunction{Play}[\cos [\pi \multsp 1800\multsp {t^2}+2\multsp \pi \multsp 500\multsp t],\{t,0,2\},\Mvariable{SampleRate}\rightarrow
8192];}
\mathGraphic{lecture23.tex_gr4.eps}

\end{CellGroup}

\end{CellGroup}

\end{CellGroup}
\begin{CellGroup}
\Section*{3 Sampling theorem}


The conditions under which a continuous-time signal can be completely represented by samples equally spaced in time are given by what is known as
the sampling theorem. This theorem is fundamentally important.

\Theorem{Let \inlineTFinmath{x(t)} be a band-limited signal with \inlineTFinmath{X(\omega )=0} for \inlineTFinmath{\VerticalSeparator \omega \VerticalSeparator
>{{\omega }_M}}. Then \inlineTFinmath{x(t)} is uniquely determined by its samples \inlineTFinmath{x(n\multsp T)}, with \inlineTFinmath{n=0,\pm 1,\pm
2,...,} iff 

			\inlineTFinmath{{{\omega }_s}\multsp >\multsp 2\multsp {{\omega }_M}}

where

			\inlineTFinmath{{{\omega }_s}=\frac{2\pi }{T}}.

Given these samples we can reconstruct \inlineTFinmath{x(t)} exactly.}


\begin{CellGroup}
\Subsubsection*{Problem 2}
(a) Given the signal \inlineTFinmath{f(t)\multsp =\multsp \frac{\sin(\pi \multsp t)}{\pi \multsp t}}, obtain the minimum sampling frequency ({\itshape
i.e.,} the Nyquist rate) and sampling period.



(b) Sample the signal at a rate 33\%{} higher than the result found in (a). Plot the resulting samples.


\end{CellGroup}

\end{CellGroup}
\begin{CellGroup}
\Section*{4 Interpolation}


To recover the original continuous-time signal following ideal sampling, the sampled signal must be passed through an ideal lowpass filter with a
cutoff at \inlineTFinmath{{{\omega }_c}=\frac{\pi }{T}}, where \inlineTFinmath{T} is the sampling period.



	\inlineTFinmath{S(\omega )\multsp =\multsp \big\{\MathBegin{MathArray}[c]{cc}
	T,&\big|\omega \big|<\frac{\pi }{T} \\
	0,&\Mvariable{else} 
  \MathEnd{MathArray}}



The impulse response of this filter is the familiar sinc function.

\dispSFinmath{\Muserfunction{sinc}[\Mvariable{t\_},\Mvariable{T\_}]:=\multsp \frac{T}{\pi \multsp t}\multsp \sin \big[\frac{\pi \multsp t}{T}\big]}
Since a filtering operation, using the ideal lowpass filter, reconstructs the original signal from its samples, this operation is a form of interpolation.
The filter is an ideal interpolator in the sense that it reconstructs the original signal exactly at all points on the real axis. The sinc function
is an interpolating function. Mathematically, interpolation is a convolution of the sampled signal with the interpolating function.

	\inlineTFinmath{x(t)\multsp =\multsp \bigg(\sum _{k=-\infty }^{\infty }x(k\multsp T)\multsp \delta (t-k\multsp T)\bigg)\multsp *\multsp \sin\big(\frac{\pi
}{T}t\big)\big/\big(\frac{\pi }{T}t\big)}



where * denotes convolution. This last expression simplifies to a superposition of scaled and shifted sinc functions:

	\inlineTFinmath{x(t)\multsp =\multsp \multsp \multsp \sum _{k=-\infty }^{\infty }x(k\multsp T)\multsp \frac{\sin\big(\frac{\pi }{T}\multsp (t-k\multsp
T)\big)}{\frac{\pi }{T}\multsp (t-k\multsp T)}}



For example, given the following samples:

	\inlineTFinmath{x(k\multsp T)\multsp =\multsp \big\{\MathBegin{MathArray}[c]{cc}
	1,&k=0 \\
	\frac{1}{2},&k=1 \\
	-1,&k=2 \\
	0,&\Mvariable{otherwise} 
  \MathEnd{MathArray}}



we obtain the following reconstruction formula  (assume \inlineTFinmath{T=1}):



\dispSFinmath{\MathBegin{MathArray}{l}
T=1;  \\
\noalign{\vspace{0.9375ex}}  \\
x[\Mvariable{t\_}]:=\multsp \Muserfunction{sinc}[t,\multsp T]+\frac{1}{2}\Muserfunction{sinc}[t-T,T]-\Muserfunction{sinc}[t-2T,T]
\MathEnd{MathArray}}
\begin{CellGroup}
\dispSFinmath{\MathBegin{MathArray}{l}
\Mfunction{Plot}\big[\big\{\Muserfunction{sinc}[t,T],\frac{1}{2}\Muserfunction{sinc}[t-T,T],-\Muserfunction{sinc}[t-2\multsp T,T],x[t]\big\},  \\
\noalign{\vspace{0.90625ex}}
\hspace{1.em} \{t,-2,6\},\Mvariable{PlotRange}->\Mvariable{All},  \\
\noalign{\vspace{0.5ex}}
\hspace{1.em} \Mvariable{PlotLabel}\rightarrow \Mvariable{Se\NTilde al\multsp reconstruida\multsp en\multsp tiempo\multsp continuo},\Mvariable{PlotStyle}->
 \\
\noalign{\vspace{0.666667ex}}
\hspace{2.em} \{\Mfunction{Hue}[0],\Mfunction{Hue}[1/3],\Mfunction{Hue}[2/3],\{\Mfunction{Thickness}[0.01],\Mfunction{GrayLevel}[0]\}\}\big];\\
\MathEnd{MathArray}}
\mathGraphic{lecture23.tex_gr5.eps}

\end{CellGroup}
The resulting signal is shown in black, while the individual contributions of the non-zero samples of the signal are shown in red, green, and blue.

As a second example consider the problem of sampling and reconstruction of the following simple signal.

\dispSFinmath{x[\Mvariable{t\_}]:=\cos \big[\frac{\pi }{2}\multsp t\big]}
\begin{CellGroup}
\dispSFinmath{\Mfunction{Plot}[x[t],\{t,-2,10\}];}
\mathGraphic{lecture23.tex_gr6.eps}

\end{CellGroup}
According to the sampling theorem, for adequate reconstruction we need to choose a sampling rate in excess of \inlineTFinmath{2} samples every 4
seconds (why?), or equivalently a sampling period less than \inlineTFinmath{T=2} seconds. Here we choose \inlineTFinmath{T=1}. The interpolation
formula calls for an infinite summation. Here we investigate the approximations to \inlineTFinmath{x(t)} in the region \inlineTFinmath{-1\leq t\leq
1}, based on finite summations, i.e. only contributions from a finite number of neighboring samples will be considered.



\dispSFinmath{\MathBegin{MathArray}{l}
T=1;  \\
\noalign{\vspace{1.0625ex}}  \\
r[\Mvariable{t\_},\Mvariable{K\_}]:=\sum _{k=-K}^{K}x[k\multsp T]\multsp \Muserfunction{sinc}[t-k\multsp T,\multsp T]
\MathEnd{MathArray}}
Here is an interpolation example using the following samples: \inlineTFinmath{\{x(-2\multsp T),\multsp x(-T),\multsp x(0),\multsp x(T),\multsp x(2\multsp
T)\}}.



\begin{CellGroup}
\dispSFinmath{r[t,2]}
\dispSFoutmath{-\frac{\sin [\pi \multsp (-2+t)]}{\pi \multsp (-2+t)}+\frac{\sin [\pi \multsp t]}{\pi \multsp t}-\frac{\sin [\pi \multsp (2+t)]}{\pi
\multsp (2+t)}}

\end{CellGroup}
\begin{CellGroup}
\dispSFinmath{\Mfunction{Plot}[x[t],\{t,-4\multsp \pi ,4\multsp \pi \},\Mvariable{PlotLabel}->\Mvariable{Se\NTilde al\multsp de\multsp tiempo\multsp
continuo}]}
\mathGraphic{lecture23.tex_gr7.eps}
\dispSFoutmath{\SkeletonIndicator \Mvariable{Graphics}\SkeletonIndicator }

\end{CellGroup}
\begin{CellGroup}
\dispSFinmath{\MathBegin{MathArray}{l}
\Muserfunction{LinePlot}\big[\Mfunction{Table}\big[\{k\multsp T,x[k\multsp T]\},\big\{k,-\frac{4}{T},\frac{5}{T}\big\}\big],  \\
\noalign{\vspace{1.07292ex}}
\hspace{1.em} \Mvariable{PlotLabel}->\Mvariable{Se\NTilde al\multsp muestreada},\Mvariable{PlotRange}->\Mvariable{All}\big]\\
\MathEnd{MathArray}}
\mathGraphic{lecture23.tex_gr8.eps}
\dispSFoutmath{\SkeletonIndicator \Mvariable{Graphics}\SkeletonIndicator }

\end{CellGroup}
\begin{CellGroup}
\dispSFinmath{\MathBegin{MathArray}{l}
\Mfunction{Plot}[r[t,8],\{t,-4\multsp \pi ,4\multsp \pi \},  \\
\noalign{\vspace{0.5ex}}
\hspace{1.em} \Mvariable{PlotLabel}->\Mvariable{Se\NTilde al\multsp reconstruida\multsp con\multsp 10\multsp muestras}]\\
\MathEnd{MathArray}}
\mathGraphic{lecture23.tex_gr9.eps}
\dispSFoutmath{\SkeletonIndicator \Mvariable{Graphics}\SkeletonIndicator }

\end{CellGroup}
Here we plot the interpolation error for three different approximations of the original signal.

\begin{CellGroup}
\dispSFinmath{\MathBegin{MathArray}{l}
\Mvariable{DisplayGraphicArrayTogether}[  \\
\noalign{\vspace{0.5ex}}
\hspace{1.em} \{\{\Mfunction{Plot}[x[t]-r[t,2],\{t,-2,2\},\Mvariable{PlotRange}\rightarrow \Mvariable{All},\Mvariable{PlotLabel}\rightarrow 2]\},\IndentingNewLine
  \\
\noalign{\vspace{0.5ex}}
\hspace{2.em} \{\Mfunction{Plot}[x[t]-r[t,5],\{t,-2,2\},\Mvariable{PlotRange}\rightarrow \Mvariable{All},\Mvariable{PlotLabel}\rightarrow 5]\},\IndentingNewLine
  \\
\noalign{\vspace{0.5ex}}
\hspace{2.em} \{\Mfunction{Plot}[x[t]-r[t,10],\{t,-2,2\},\Mvariable{PlotRange}\rightarrow \Mvariable{All},\Mvariable{PlotLabel}\rightarrow 10]\}\}]\\
\MathEnd{MathArray}}
\begin{CellGroup}
\mathGraphic{lecture23.tex_gr10.eps}
\mathGraphic{lecture23.tex_gr11.eps}
\mathGraphic{lecture23.tex_gr12.eps}

\end{CellGroup}
\dispSFoutmath{\MathBegin{MathArray}{l}
\Mvariable{DisplayGraphicArrayTogether}[  \\
\noalign{\vspace{0.5ex}}
\hspace{1.em} \{\{\SkeletonIndicator \Mvariable{Graphics}\SkeletonIndicator \},\{\SkeletonIndicator \Mvariable{Graphics}\SkeletonIndicator \},\{\SkeletonIndicator
\Mvariable{Graphics}\SkeletonIndicator \}\}]\\
\MathEnd{MathArray}}

\end{CellGroup}
Here we measure the power of the approximation error  defined as the integral squared error over one period of the signal \inlineTFinmath{x(t)}.
Note that due to the symmetry of the problem, we can simplify the error calculation by evaluating over one quarter of the period only.



\begin{CellGroup}
\dispSFinmath{(x[t]-r[t,2])\RawWedge 2}
\dispSFoutmath{{{\Big(\cos \big[\frac{\pi \multsp t}{2}\big]+\frac{\sin [\pi \multsp (-2+t)]}{\pi \multsp (-2+t)}-\frac{\sin [\pi \multsp t]}{\pi
\multsp t}+\frac{\sin [\pi \multsp (2+t)]}{\pi \multsp (2+t)}\Big)}^2}}

\end{CellGroup}
\begin{CellGroup}
\dispSFinmath{\Mvariable{NP}=\frac{1}{2}\Mfunction{NIntegrate}[\%,\multsp \{t,0,2\}]\multsp }
\dispSFoutmath{0.000736946}

\end{CellGroup}
The signal power is 

\dispSFinmath{\Mvariable{SP}=0.5}
The signal-to-noise ratio (in dB) is then given by 

\begin{CellGroup}
\dispSFinmath{10\multsp \log \big[10,\multsp \frac{\Mvariable{SP}}{\Mvariable{NP}}\big]}
\dispSFoutmath{28.3153}

\end{CellGroup}
\begin{CellGroup}
\Subsubsection*{Problem 3}
Measure the approximation error power and the signal-to-noise ratio (SNR) as a function of K, the number of samples used in the signal reconstruction.


\end{CellGroup}

\end{CellGroup}
\Section*{PROBLEM SOLUTIONS}
\Section*{HOMEWORK PROBLEMS}

\end{document}
